% !TeX root = ../main.tex
\documentclass[../main.tex]{subfiles}

\begin{document}
\ifSubfilesClassLoaded{\standalonehead}{}

\lecture{1}{2 сентября 2026}{Амортизированный анализ}

Оценка «худший случай одной операции» бывает слишком грубой: дорогая операция может
быть редкой, и платить за неё приходится не на каждом шаге. Амортизированный анализ
оценивает не отдельную операцию, а стоимость всей последовательности.

\subsection{Амортизированная стоимость}

\begin{Definition}{Амортизированная стоимость}{amort}
  Пусть структура данных обрабатывает последовательность операций
  $\mathrm{op}_1,\dots,\mathrm{op}_m$ с фактическими стоимостями $c_1,\dots,c_m$.
  Число $\Amort$ называется \term{амортизированной стоимостью} операции, если
  для любой допустимой последовательности длины $m$
  \[
    \sum_{i=1}^{m} c_i \;\le\; m\cdot\Amort .
  \]
\end{Definition}

\begin{Remark}{}{avg}
  Это не усреднение по вероятности: случайность нигде не используется, оценка
  обязана выполняться для \emph{худшей} последовательности операций. Отдельная
  операция при этом может стоить сколь угодно дорого.
\end{Remark}

\subsection{Динамический массив}

Массив хранит $n$ элементов в буфере ёмкости $N$. Вставка в конец стоит $1$, если
$n < N$; иначе выделяется буфер вдвое большей ёмкости, старое содержимое
копируется, и операция стоит $n+1$.

\begin{figure}[H]
  \centering
  \begin{tikzpicture}
    \begin{axis}[notestyle, height=4.6cm, width=0.86\linewidth,
      xlabel={номер операции $i$}, ylabel={стоимость $c_i$},
      xmin=0.2, xmax=16.8, ymin=0, ymax=10,
      xtick={1,4,8,12,16}, ytick={0,3,6,9},
      grid=none]
      \addplot[ybar, bar width=6pt, fill=figA!25, draw=figA] coordinates {
        (1,1) (2,2) (3,3) (4,1) (5,5) (6,1) (7,1) (8,1)
        (9,9) (10,1) (11,1) (12,1) (13,1) (14,1) (15,1) (16,1)};
      \addplot[sharp plot, figC, dashed, line width=0.9pt, mark=none]
        coordinates {(0.2,3) (16.8,3)};
      \node[lblS, anchor=west, text=figC] at (axis cs:10.2,3.9) {$\Amort = 3$};
    \end{axis}
  \end{tikzpicture}
  \caption{Фактическая стоимость вставок в динамический массив. Пики приходятся на
    шаги, где ёмкость удваивается; их вклад покрывается постоянной амортизированной
    стоимостью.}
  \label{fig:pushback}
\end{figure}

\begin{Theorem}{Динамический массив}{dynarray}
  Последовательность из $n$ вставок в конец, начатая с пустого массива ёмкости~$1$,
  выполняется суммарно не более чем за $3n$ элементарных операций. Амортизированная
  стоимость вставки есть $\bigO(1)$.
\end{Theorem}

\begin{proof}
  Разделим стоимость на две части: собственно запись элемента и копирование при
  расширении.

  Записей ровно $n$, каждая стоит $1$: вклад $n$.

  Копирование происходит, когда текущий размер равен ёмкости, то есть при размерах
  $1, 2, 4, \dots, 2^{t}$, где $2^{t} \le n$. Копируется столько элементов, каков
  размер, поэтому суммарный вклад копирований равен
  \dop[сумма геометрической прогрессии]
  \[
    \sum_{j=0}^{t} 2^{j} \;=\; 2^{t+1} - 1 \;<\; 2\cdot 2^{t} \;\le\; 2n .
  \]
  В сумме получаем строго меньше $3n$, что и требовалось. Деля на $n$, получаем
  $\Amort < 3$.
\end{proof}

\begin{algorithm}[H]
  \caption{Вставка в конец динамического массива}
  \label{alg:pushback}
  \KwIn{массив $A$ ёмкости $N$, размер $n$, элемент $x$}
  \If{$n = N$}{
    выделить буфер $B$ ёмкости $2N$\;
    \For{$i \gets 0$ \KwTo $n-1$}{$B[i] \gets A[i]$\;}
    $A \gets B$; \quad $N \gets 2N$\;
  }
  $A[n] \gets x$; \quad $n \gets n+1$\;
\end{algorithm}

\subsection{Метод потенциалов}

Агрегатный подсчёт из теоремы~\ref{thm:dynarray} работает, пока последовательность
операций однородна. Метод потенциалов снимает это ограничение: он приписывает
состоянию структуры число --- «запас оплаченной вперёд работы».

\begin{Definition}{Потенциал}{potential}
  Пусть $D_i$ --- состояние структуры после $i$-й операции. Функция
  $\Pot\colon \{D_i\} \to \RR$ называется \term{потенциалом}, если
  $\Pot(D_0) = 0$ и $\Pot(D_i) \ge 0$ для всех $i$. Амортизированной стоимостью
  $i$-й операции относительно $\Pot$ называют
  \[
    \Amort_i \;=\; c_i + \Pot(D_i) - \Pot(D_{i-1}).
  \]
\end{Definition}

\begin{Lemma}{}{potbound}
  Для любого потенциала $\Pot$ и любой последовательности из $m$ операций
  \[
    \sum_{i=1}^{m} c_i \;\le\; \sum_{i=1}^{m} \Amort_i .
  \]
\end{Lemma}

\begin{proof}
  Сумма телескопируется:
  \[
    \sum_{i=1}^{m}\Amort_i
      = \sum_{i=1}^{m} c_i + \Pot(D_m) - \Pot(D_0)
      = \sum_{i=1}^{m} c_i + \Pot(D_m) \;\ge\; \sum_{i=1}^{m} c_i,
  \]
  поскольку $\Pot(D_0) = 0$ и $\Pot(D_m) \ge 0$.
\end{proof}

Смысл леммы~\ref{lem:potbound} в том, что верхняя оценка на каждое $\Amort_i$ по
отдельности автоматически даёт оценку на всю последовательность.

\begin{Example}{Двоичный счётчик}{counter}
  Счётчик хранится как массив битов $b_0, b_1, \dots$; операция
  $\textsc{Increment}$ переворачивает младшие единицы в нули и следующий ноль
  в единицу. Стоимость --- число перевёрнутых битов --- меняется от $1$ до
  $\Theta(\log v)$.

  \vspace{4pt}
  \centering
  \begin{tikzpicture}[x=0.95cm, y=0.8cm]
    \node[lblS, anchor=east] at (-0.55, 0.95) {$v$};
    \foreach \b in {0,1,2,3}{\node[lblS] at (3-\b, 0.95) {$b_\b$};}
    \node[lblS, anchor=west] at (3.95, 0.95) {$c_i$};
    \foreach \v/\k in {0/{}, 1/1, 2/2, 3/1, 4/3, 5/1, 6/2, 7/1, 8/4}{
      \node[lblS, anchor=east] at (-0.55, -\v) {$\v$};
      \foreach \b in {0,1,2,3}{
        \pgfmathtruncatemacro{\bit}{mod(floor(\v/(2^\b)),2)}
        \pgfmathtruncatemacro{\fl}{(\v>0) && (mod(\v-1,2^\b)==(2^\b-1)) ? 1 : 0}
        \ifnum\fl=1
          \node[cell, fill=figA!15, draw=figA, text=figA] at (3-\b, -\v) {\bit};
        \else
          \node[cell, text=inkSoft] at (3-\b, -\v) {\bit};
        \fi
      }
      \node[lblS, anchor=west] at (3.95, -\v) {\k};
    }
  \end{tikzpicture}
  \vspace{2pt}

  \raggedright\footnotesize\color{inkSoft}
  Синим отмечены биты, перевёрнутые на текущем шаге. Дорогие шаги
  ($v = 4$, $v = 8$) наступают тем реже, чем они дороже.
\end{Example}

\begin{Statement}{}{counterbound}
  Амортизированная стоимость $\textsc{Increment}$ равна $2$: $n$~инкрементов
  переворачивают менее $2n$ битов.
\end{Statement}

\begin{proof}
  Возьмём в качестве потенциала число единиц в записи счётчика:
  $\Pot(D_i) = \#\{b_j = 1\}$. Условия определения~\ref{def:potential} выполнены:
  $\Pot(D_0) = 0$ и $\Pot \ge 0$.

  Пусть на $i$-м шаге сброшено $k$ единиц; тогда переворачивается ещё один ноль,
  так что $c_i = k + 1$, а число единиц меняется на $\Pot(D_i) - \Pot(D_{i-1}) = 1 - k$.
  Следовательно,
  \[
    \Amort_i = (k+1) + (1-k) = 2 .
  \]
  По лемме~\ref{lem:potbound} суммарное число переворотов не превосходит $2n$.
\end{proof}

\begin{Addition}[выбор коэффициента роста]
  На лекции разбирался только рост в два раза. Тот же подсчёт проходит для любого
  коэффициента $\alpha > 1$: копирования дают
  $\sum_j \alpha^{\,j} = \bigO(\alpha^{\,t}) = \bigO(n)$, поэтому
  $\Amort = \bigO\!\left(\frac{\alpha}{\alpha - 1}\right)$ --- константа, зависящая
  только от~$\alpha$. Практическое значение $\alpha = 1{,}5$ (используется, например,
  в некоторых реализациях \eng{std::vector}) даёт большую константу, но позволяет
  переиспользовать ранее освобождённую память.

  А вот рост \emph{на константу} $\Delta$ ломает оценку: расширения происходят на
  шагах $\Delta, 2\Delta, \dots$, и суммарная стоимость копирований равна
  $\Theta(n^2/\Delta)$, то есть $\Amort = \Theta(n)$.
\end{Addition}

\end{document}
